%% ma: L11-B26-0003
%% lop: 11
%% chuong: 7
%% bai: 26
%% dang: 11.26.3
%% loai: TLN
%% muc: VD
%% dap_an: "0,24"
%% nguon: "(NBV)-ĐỀ SỐ 2-MỖI NGÀY 1 ĐỀ THI 2025.pdf (D02, MỖI NGÀY 1 ĐỀ - Đề số 2), Phần III Câu 1 (THPT Văn Giang - Hưng Yên 2025)"
%% kien_thuc: "Bài 26 (khoảng cách giữa hai đường thẳng chéo nhau, quy về khoảng cách từ điểm đến mặt phẳng), Bài 24 (góc giữa đường thẳng và mặt phẳng), Bài 23 (đường thẳng vuông góc với mặt phẳng)"
%% ghi_chu: "Đề không có hình (hình trong lời giải gốc không cần). Đáp án gốc 0,24 đúng (kiểm bằng sympy: d=√10/5 a)"
%% trang_thai: chinh-thuc
\begin{ex}
Cho hình chóp $S.ABCD$ có đáy $ABCD$ là hình chữ nhật với $AB=a$, $AD=2a$. Hình chiếu vuông góc của $S$ trên mặt phẳng đáy là trung điểm $H$ của $AD$, góc giữa $SB$ và mặt phẳng $(ABCD)$ là $45^\circ$. Khoảng cách giữa hai đường thẳng $SD$ và $BH$ theo $a$ được kết quả là $ma$. Khi đó giá trị $\dfrac35m^2$ bằng bao nhiêu?
\shortans{0,24}
\loigiai{$SH\perp(ABCD)$ nên $\widehat{SBH}=45^\circ$, $HB=\sqrt{AB^2+AH^2}=\sqrt2\,a$, suy ra $SH=HB=\sqrt2\,a$.
Gọi $M$ là trung điểm $BC$. $HD\parallel BM$, $HD=BM=a$ nên $HBMD$ là hình bình hành, $DM\parallel HB$, do đó $HB\parallel(SDM)$ và $d(HB,SD)=d(H,(SDM))$.
$HDCM$ là hình vuông tâm $O$ nên $HC\perp DM$; lại có $SH\perp DM$ nên $DM\perp(SHC)$. Kẻ $HI\perp SO$ ($I\in SO$) thì $HI\perp DM$ nên $HI\perp(SDM)$, suy ra $d(H,(SDM))=HI$.
Có $HO=\dfrac12HC=\dfrac{\sqrt2}{2}a$, $HI=\dfrac{SH\cdot HO}{\sqrt{SH^2+HO^2}}=\dfrac{a^2}{\sqrt{\frac52}\,a}=\dfrac{\sqrt{10}}{5}a$.
Vậy $m=\dfrac{\sqrt{10}}5$, $\dfrac35m^2=\dfrac35\cdot\dfrac{10}{25}=0{,}24$.}
\end{ex}
